\section{Removing a determinant factor}\label{sec:reduced}

The polynomial $p$ makes the positive-map construction explicit through
a $4\times4$ determinant. We now remove a factor that is unnecessary for
its hyperbolicity cone. A larger block determinant shows directly that
the quotient remains polynomial at singular $X$.

\begin{proof}[Proof of Corollary~\ref{cor:reduced}]
For the standard coordinate vectors $e_i$ of $\R^3$, put
\[
 K_i=\begin{pmatrix}0&e_i^\top\\-e_i&0\end{pmatrix}\in\R^{4\times4},
 \qquad C=\begin{pmatrix}K_1&K_2&K_3\end{pmatrix}\in\R^{4\times12}.
\]
Here $U\otimes V$ denotes the block matrix with blocks $U_{ij}V$.
Direct block multiplication in~\eqref{eq:phi} gives, for $T\in\Sym^4$,
\[
 \Phi_y(T)=\sum_{i,j=1}^3 Q(y)_{ij}K_iTK_j^\top
          =C\bigl(Q(y)\otimes T\bigr)C^\top.
\]
Define the polynomial
\begin{equation}\label{eq:reduced-polynomial}
 q(X,Z,y)=\det\begin{pmatrix}
 I_3\otimes X&C^\top\\
 C\bigl(Q(y)\otimes I_4\bigr)&Z
 \end{pmatrix}.
\end{equation}
For invertible $X$, taking a Schur complement yields
\[
 q(X,Z,y)=(\det X)^3\det\bigl(Z-\Phi_y(X^{-1})\bigr).
\]
Comparison with~\eqref{eq:inverse-formula} proves the polynomial identity
\begin{equation}\label{eq:reduced-factor}
 p(X,Z,y)=(\det X)q(X,Z,y)
\end{equation}
first for invertible $X$ and then everywhere by continuity.
Since $p$ and $\det X$ are homogeneous of degrees $20$ and $4$, this
identity makes $q$ homogeneous of degree $16$; also $q(e)=1$.
For each real $(X,Z,y)$, the polynomial $q(te-(X,Z,y))$ is a factor of
the real-rooted polynomial $p(te-(X,Z,y))$. Hence $q$ is hyperbolic
with respect to $e$.

Write $K_q=\Lambda_+(q,e)$. The factorization~\eqref{eq:reduced-factor}
gives
\begin{equation}\label{eq:cone-intersection}
 K=\{(X,Z,y):X\succeq0\}\cap K_q,
\end{equation}
because the roots contributed by $\det(tI_4-X)$ are the eigenvalues of
$X$. It remains to show that $K_q$ already forces $X\succeq0$.
The polynomial $q$ is even in $y$, so $(X,Z,y)\in K_q$ implies
$(X,Z,-y)\in K_q$. By convexity of hyperbolicity
cones~\cite{garding-1959}, their midpoint $(X,Z,0)$ also belongs to $K_q$.
But~\eqref{eq:reduced-polynomial} gives
\[
 q(X,Z,0)=(\det X)^3\det Z,
\]
whose closed hyperbolicity cone in this slice is exactly
$\{(X,Z,0):X\succeq0,\ Z\succeq0\}$.
Thus $X\succeq0$ throughout $K_q$, and~\eqref{eq:cone-intersection}
implies $K_q=K$. Theorem~\ref{thm:main} now supplies nonspectrahedrality.
\end{proof}
